\documentclass[11pt]{article}

% Change "review" to "final" to generate the final (sometimes called camera-ready) version.
% Change to "preprint" to generate a non-anonymous version with page numbers.
\usepackage[review]{acl}

% Standard package includes
\usepackage{times}
\usepackage{latexsym}
\usepackage[T1]{fontenc}
\usepackage[utf8]{inputenc}
\usepackage{microtype}
\usepackage{inconsolata}
\usepackage{graphicx}
\usepackage{amsmath}
\usepackage{amssymb}

\title{Anything2PBL: Hierarchical Constrained Generation with \\ Learner-Aware Scaffolding for Project-Based Learning}

\author{First Author \\
  Affiliation / Address line 1 \\
  Affiliation / Address line 2 \\
  \texttt{email@domain} \\\And
  Second Author \\
  Affiliation / Address line 1 \\
  Affiliation / Address line 2 \\
  \texttt{email@domain} \\}

\begin{document}
\maketitle

\begin{abstract}
Project-Based Learning (PBL) is an effective approach for developing practical competencies through authentic project work.
However, creating high-quality training projects remains costly, as it requires not only engineering completeness but also learner-appropriate pedagogical design---in particular, deciding what constitutes a productive challenge versus tedious grunt work depends on the target learner rather than the project itself.
Existing approaches either focus on code generation without pedagogical consideration, or assist the learning process without producing a complete, runnable project.
We propose \textbf{Anything2PBL} (A2P), a hierarchical constrained generation framework that transforms arbitrary inputs into pedagogically-calibrated project repositories through a five-stage human-in-the-loop pipeline.
Experiments on a cross-disciplinary benchmark show that A2P consistently outperforms direct-prompting and chain-of-thought baselines, and that instructor calibration at the scaffolding stage yields disproportionately large gains.
\end{abstract}


\bibliography{custom}

\end{document}
